\documentclass[12pt, a4paper]{article}

\usepackage[T2A]{fontenc}
\usepackage[utf8]{inputenc}
\usepackage[bulgarian]{babel}
\usepackage{geometry}
\geometry{a4paper, margin=1in}
\usepackage{tabularx}
\usepackage{xltabular}
\usepackage{ragged2e}
\usepackage{booktabs}
\usepackage{hyperref}

\title{Задание 4: Персони и потребителски сценарии}
\author{}
\date{}

\begin{document}

\maketitle

\noindent\textbf{Дисциплина:} Проектиране на човеко-машинен интерфейс 2025-2026

\bigskip

\noindent\textbf{Участници в проекта:}
\smallskip

\noindent
\begin{tabularx}{\textwidth}{|l|X|l|X|l|}
\hline
\textbf{№} & \textbf{Име и фамилия} & \textbf{Факултетен №} & \textbf{Специалност} & \textbf{Курс} \\ \hline
1 & Виктор Койчев & 3MI0600164 & Софтуерно инженерство & 3-ти \\ \hline
2 & Емилиан Янев & 0MI0600273 & Софтуерно инженерство & 3-ти \\ \hline
3 & Жан Георгиев & 8MI0600398 & Софтуерно инженерство & 3-ти \\ \hline
4 & Даниел Николов & 0MI0600410 & Софтуерно инженерство & 3-ти \\ \hline
\end{tabularx}

\bigskip

\noindent\textbf{Име на група:}\\
HCI\_2026\_group\_3MI0600164\_0MI0600273\_8MI0600398\_0MI0600410

\bigskip

\noindent\textbf{Име на проекта:}\\
Система за управление и мониторинг на роботизиран склад за фармацевтични продукти (PharmaBot Logistics)

\bigskip
\rule{\textwidth}{0.4pt}

\section*{1. Анализ на потребностите чрез Персони (Personas)}

За нуждите на проекта разработихме два основни потребителски профила (Персони), които представляват крайните потребители на двата различни интерфейса на системата: мобилното приложение на товарната рампа и десктоп дашборда в контролната зала.

\bigskip
\noindent\textbf{Персона 1: Шофьор на камион (Мобилен интерфейс)}

\begin{itemize}
    \item \textbf{Име:} Иван Димитров
    \item \textbf{Възраст:} 48 години
    \item \textbf{Професия:} Шофьор в логистична компания
    \item \textbf{Технически умения:} Ниски до средни. Използва смартфон предимно за базови комуникации (Viber, навигация). Често се затруднява със сложни корпоративни софтуери.
    \item \textbf{Работна среда:} На открито, често на пряка слънчева светлина или дъжд. Работи с предпазни ръкавици. Бърза, защото графикът му е натоварен.
    \item \textbf{Цели (Goals):}
    \begin{itemize}
        \item Да натовари стоката възможно най-бързо и да потегли.
        \item Да има недвусмислено доказателство, че е приел палета в изрядно състояние, за да не носи вина при евентуален проблем.
    \end{itemize}
    \item \textbf{Проблеми и Фрустрации (Pain points):}
    \begin{itemize}
        \item Губи средно по 20 минути на курс в чакане на хартиени протоколи и подписи.
        \item Дразни се от приложения с малки бутони, които не може да натисне с ръкавици.
        \item Трудно чете от екрана на таблета си през лятото заради отблясъците от слънцето.
    \end{itemize}
\end{itemize}

\bigskip
\noindent\textbf{Персона 2: Складов оператор (Десктоп Дашборд)}

\begin{itemize}
    \item \textbf{Име:} Елена Георгиева
    \item \textbf{Възраст:} 34 години
    \item \textbf{Професия:} Оператор „Складова логистика и мониторинг“
    \item \textbf{Технически умения:} Високи. Работи на компютър ежедневно, свикнала е да следи множество данни на няколко монитора едновременно.
    \item \textbf{Работна среда:} Климатизирана контролна зала в склада. Средата е спокойна, докато не възникне инцидент (разлив, блокиран път, технически отказ).
    \item \textbf{Цели (Goals):}
    \begin{itemize}
        \item Да осигури безпроблемното движение на флотилия от над 50 робота едновременно.
        \item Строго да следи за спазването на температурния режим (Cold Chain) на чувствителните медикаменти.
    \end{itemize}
    \item \textbf{Проблеми и Фрустрации (Pain points):}
    \begin{itemize}
        \item При инцидент в склада ѝ отнема твърде много време да въведе ръчно забранените координати в старата система, което води до "задръствания" от роботи.
        \item Получава твърде много известия едновременно и понякога трудно различава критичните температурни аларми от стандартните нотификации за изтощена батерия.
    \end{itemize}
\end{itemize}

\bigskip
\rule{\textwidth}{0.4pt}

\section*{2. Потребителски сценарии (User Scenarios)}

Въз основа на изградените персони, дефинираме следните сценарии, които демонстрират как системата \textit{PharmaBot Logistics} решава техните проблеми.

\bigskip
\noindent\textbf{Сценарий 1: Дигитално предаване на открито (Иван)}

\smallskip
\noindent\textbf{Контекст:} Юли месец, 14:00 часа. Слънцето грее ярко директно в товарна рампа №4. Иван пристига с камиона си, за да натовари партида скъпи ваксини.

\smallskip
\noindent\textbf{Действие:} Иван слиза от камиона с работните си ръкавици и отваря приложението PharmaBot на служебния таблет. Заради силното слънце той не вижда добре екрана, затова веднага натиска бутона „Контраст“ (Dark Mode), който прави интерфейса черен с ярки бели букви – проблемът с четимостта е решен. 

До него пристига автономният робот (BOT-042) с палета. Иван вижда, че само първата стъпка („Сканирай QR код“) е активна – големият син бутон е лесен за натискане дори с ръкавици. Той сканира кода от екрана на робота. Системата моментално отключва втората стъпка. Иван натоварва палета, натиска бутона „Снимай палета“, прави снимка за доказателство и потвърждава експедицията.

\smallskip
\noindent\textbf{Резултат:} Целият процес му отнема 45 секунди. Той потегля веднага, без да чака администратор за хартиени подписи, уверен, че дигиталното доказателство го пази от отговорност.

\bigskip
\noindent\textbf{Сценарий 2: Кризисна реакция и спасяване на „Студената верига“ (Елена)}

\smallskip
\noindent\textbf{Контекст:} Елена е на смяна в контролната зала и наблюдава Сектор А през 2D интерактивната карта. Всички 50 робота се движат нормално.

\smallskip
\noindent\textbf{Действие:} Изведнъж работник с мотокар изпуска палет по средата на главната пътека, блокирайки трафика. Елена веднага цъква бутона „Режим: Очертай зона“ и с мишката начертава червен квадрат върху засегнатия участък на картата. Системата автоматично и на секундата преизчислява маршрутите на всички приближаващи роботи, предотвратявайки сблъсъци. 

В същия момент екранът ѝ се затъмнява и изскача критичен червен модален прозорец: робот BOT-019, пренасящ температуро-чувствителен инсулин, отчита повреда в охладителната система (температура 12.5°C при позволени 8°C). Без да се налага да отваря допълнителни менюта, Елена кликва директно върху бутона „Пренасочи към Карантина“ в самия прозорец с алармата.

\smallskip
\noindent\textbf{Резултат:} Елена успява да овладее физически инцидент (разлив) и хардуерен дефект (охлаждане) в рамките на 10 секунди, спасявайки скъпоструваща пратка и предотвратявайки логистичен хаос, благодарение на интуитивния визуален контрол.

\end{document}